\documentclass[11pt]{article}

\usepackage[margin=1in]{geometry}
\usepackage{amsmath, amssymb}
\usepackage[T1]{fontenc}
\usepackage[utf8]{inputenc}
\usepackage{lmodern}
\usepackage{enumitem}
\usepackage{hyperref}

\setlist[itemize]{topsep=4pt, itemsep=2pt, parsep=0pt}

\title{6.8610 Spring 2026 Lecture 3 Notes: Recurrent Neural Networks, Seq2Seq, and Attention}
\author{Junru Ren with help from Codex}
\date{February 10, 2026}

\begin{document}
\maketitle

\section{Recurrent Neural Networks, Seq2Seq, and Attention}

\subsection{Motivation}
The lecture introduces \textbf{recurrent neural networks (RNNs)} as the first neural architecture in the course that is explicitly designed for sequential data. The immediate motivation is that feed-forward models with fixed windows are poorly matched to language: they do not retain state across predictions, they scale poorly with larger contexts, and they struggle with long-range dependencies.

\begin{itemize}
    \item Language is fundamentally sequential: characters form words, words form sentences, and sentences form larger discourse units.
    \item Many important linguistic phenomena depend on long-range context. For example, pronoun choice in ``Emily \dots proud of her'' versus ``Miguel \dots proud of him'' depends on a word many tokens earlier.
    \item Feed-forward models can use nonlinearities, but they still rely on fixed context windows and do not maintain memory from one prediction to the next.
    \item RNNs address this by carrying forward a hidden state that summarizes prior context.
    \item Once RNNs are available, they can be extended from next-token prediction to more general sequence modeling, including variable-length input/output tasks such as machine translation.
    \item This leads naturally to \textbf{sequence-to-sequence (seq2seq)} models and then to \textbf{attention}, which solves the bottleneck of compressing the entire source sentence into a single vector.
\end{itemize}

\subsection{Key Concepts}
\textbf{Definition: Word Token vs.\ Word Type}

A \textbf{word token} is a specific occurrence of a word in text, while a \textbf{word type} is the unique lexical form. For example, in ``I ran and ran and ran'' there are 6 word tokens but only 3 word types.

\begin{itemize}
    \item \textbf{Type-based vs.\ token-based representations:}
    \begin{itemize}
        \item Word2vec-style embeddings are type-based: every occurrence of a word like ``bank'' shares one representation.
        \item RNN hidden states are token-based or contextualized: each occurrence of a word gets a context-specific representation.
    \end{itemize}
    \item \textbf{Feed-forward language models are limited:}
    \begin{itemize}
        \item fixed input length,
        \item no state carried across predictions,
        \item parameter count grows with window size,
        \item no natural parameter sharing across positions,
        \item effectively bag-of-words unless position is encoded explicitly.
    \end{itemize}
    \item \textbf{RNN core idea:} use the previous hidden state together with the current input to compute the new hidden state.
    \item \textbf{Hidden state as meaning summary:} the hidden representation is meant to encode the meaning of the sequence seen so far.
    \item \textbf{Teacher forcing:} during training, the model receives the true previous token rather than its own predicted token.
    \item \textbf{Backpropagation through time (BPTT):} training an RNN requires propagating gradients backward through the unrolled sequence.
    \item \textbf{Truncated BPTT:} for efficiency, backpropagation may be limited to a finite recent window while the forward pass still carries recurrent state.
    \item \textbf{Sequence-to-sequence (Seq2Seq):} use one RNN as an encoder for the source sequence and another as a decoder for the target sequence.
    \item \textbf{Attention:} let each decoder step compute a weighted summary of all encoder hidden states, rather than relying only on the final encoder state.
\end{itemize}

\subsection{Core Models and Equations}
\textbf{Feed-forward language model recap.} For a fixed context representation \(X\),
\[
h = \phi(WX + b_h)
\]
\[
\hat{y} = \operatorname{softmax}(Uh + b_y)
\]
This allows nonlinear combination of a fixed window, but it does not maintain sequential state.

\textbf{Basic RNN update.} For input embedding \(x_t\), hidden state \(h_t\), and output distribution \(\hat{y}_t\),
\[
h_t = f(Wx_t + Vh_{t-1} + b_h)
\]
\[
\hat{y}_t = \operatorname{softmax}(Uh_t + b_y)
\]
with initial state typically set to
\[
h_0 = \mathbf{0}.
\]

\textbf{Cross-entropy (CE) loss at each timestep.}
\[
\operatorname{CE}(y_t,\hat{y}_t) = - \sum_{j \in V} y_{t,j}\log \hat{y}_{t,j}
\]
where \(V\) is the vocabulary and \(y_t\) is the one-hot gold distribution.

\textbf{Sequence loss.}
\[
\mathcal{L} = \sum_{t=1}^{T} \operatorname{CE}(y_t,\hat{y}_t)
\]

\textbf{Backpropagation through time intuition.} Because \(h_t\) depends recursively on earlier states, gradients to earlier recurrent weights involve long chains of multiplications. Abstractly,
\[
\frac{\partial \mathcal{L}_T}{\partial h_k}
= \frac{\partial \mathcal{L}_T}{\partial h_T}
\prod_{t=k+1}^{T}
\frac{\partial h_t}{\partial h_{t-1}}
\]
which explains why gradients may vanish or explode.

\textbf{Gradient clipping.} If gradient norm exceeds threshold \(\tau\), rescale:
\[
g \leftarrow g \cdot \min\left(1, \frac{\tau}{\lVert g \rVert}\right)
\]
This preserves direction while limiting magnitude.

\textbf{Seq2Seq encoder-decoder formulation.}
\begin{itemize}
    \item The encoder consumes source tokens \(x_1,\dots,x_N\) and produces hidden states \(h^{enc}_1,\dots,h^{enc}_N\).
    \item In the basic seq2seq model, the final encoder hidden state initializes the decoder:
    \[
    h^{dec}_0 = h^{enc}_N
    \]
    \item The decoder then produces target tokens \(y_1,\dots,y_M\) autoregressively:
    \[
    h^{dec}_t = f(W_y y_{t-1} + V h^{dec}_{t-1} + b)
    \]
    \[
    \hat{y}_t = \operatorname{softmax}(U h^{dec}_t + b_y)
    \]
\end{itemize}

\textbf{Attention scores.} For decoder state \(h^{dec}_t\) and encoder state \(h^{enc}_i\), compute a raw compatibility score
\[
e_{t,i} = \operatorname{score}(h^{dec}_t, h^{enc}_i)
\]
where the score function can be a feed-forward network, a dot product, or another learned similarity function.

\textbf{Attention weights.}
\[
a_{t,i} = \frac{\exp(e_{t,i})}{\sum_{j=1}^{N}\exp(e_{t,j})}
\]
These weights form a distribution over source positions for each decoder step \(t\).

\textbf{Context vector.}
\[
c_t = \sum_{i=1}^{N} a_{t,i} h^{enc}_i
\]
This is a weighted summary of the encoder hidden states.

\textbf{Decoder with attention.} The output at decoder step \(t\) is based on both the current decoder state and the context vector:
\[
\hat{y}_t = g(h^{dec}_t, c_t)
\]
where \(g\) may be implemented by concatenation followed by an affine transformation and softmax.

\textbf{Historical machine-translation framing.} Before neural seq2seq systems, translation was often written as
\[
p(y \mid x) \propto p(x \mid y)p(y)
\]
with a separate translation model and target-language model. Neural machine translation replaced this pipeline with a single end-to-end encoder-decoder model.

\subsection{Algorithms / Architectures}
\textbf{1. Feed-forward next-token prediction}
\begin{itemize}
    \item Choose a fixed-size context window.
    \item Convert the window into embeddings or a fixed feature vector.
    \item Pass it through one or more feed-forward layers.
    \item Predict the next token with a softmax.
    \item Repeat with a sliding window over the corpus.
\end{itemize}

\textbf{2. RNN forward pass}
\begin{itemize}
    \item Initialize \(h_0\), often to zeros.
    \item For each token \(x_t\), compute \(h_t\) from the current input and previous hidden state.
    \item Produce \(\hat{y}_t\) from \(h_t\).
    \item Interpret \(h_t\) as the contextualized representation of the sequence prefix up to \(t\).
\end{itemize}

\textbf{3. RNN training}
\begin{itemize}
    \item Feed the gold token sequence one token at a time.
    \item Compute a prediction at each timestep.
    \item Evaluate cross-entropy against the gold next token.
    \item Sum losses across timesteps.
    \item Run BPTT to update \(W\), \(V\), \(U\), and optionally the input embedding matrix.
    \item Use teacher forcing during training: feed the true previous token, not the model's sampled output.
\end{itemize}

\textbf{4. Truncated BPTT}
\begin{itemize}
    \item Keep the normal recurrent forward pass.
    \item Backpropagate only through the most recent \(T\) timesteps.
    \item Update parameters periodically, for example every \(T\) tokens.
    \item This limits compute without reducing the forward-pass memory to an \(n\)-gram window.
\end{itemize}

\textbf{5. RNN generation}
\begin{itemize}
    \item Begin with a special start token such as \texttt{<START>} or the beginning-of-sequence token \texttt{<BOS>}.
    \item Produce a distribution over the next token from \(\hat{y}_1\).
    \item Sample or choose the next token.
    \item Feed that generated token back in as the next input.
    \item Continue until generating an end-of-sequence token such as \texttt{<EOS>}.
\end{itemize}

\textbf{6. Seq2Seq translation}
\begin{itemize}
    \item Run the encoder over the source sentence.
    \item Use the final encoder state to initialize the decoder.
    \item Start the decoder with a start symbol.
    \item During training, input the gold previous target token to the decoder.
    \item During testing, input the decoder's own previous prediction.
    \item Stop when the decoder emits a stop token.
\end{itemize}

\textbf{7. Attention-augmented decoding}
\begin{itemize}
    \item For each decoder step, score the current decoder state against every encoder state.
    \item Softmax the scores to produce attention weights.
    \item Compute a weighted sum of encoder states to obtain \(c_t\).
    \item Combine \(c_t\) with the decoder hidden state to predict the next target token.
    \item Repeat for each output position, producing a different context vector at each decoder step.
\end{itemize}

\subsection{Important Intuitions from the Professor}
\begin{itemize}
    \item \textbf{RNNs are the first truly contextualized embedding model in the course.} The hidden states are not just intermediates; they are the contextualized token representations.
    \item \textbf{Using bytes directly is possible but often an unnecessary uphill battle.} The professor emphasized that byte-level modeling makes meaning harder to recover because single bytes are not natural semantic units.
    \item \textbf{Teacher forcing is not a minor implementation detail.} During training, feeding model predictions back into the decoder too early would mean learning from garbage outputs instead of the correct target sequence.
    \item \textbf{Truncated BPTT does not mean truncated forward memory.} The model still carries recurrent state forward through long sequences; only the gradient computation is cut off.
    \item \textbf{Long short-term memory (LSTM) networks were a pragmatic architectural repair.} They preserve the high-level RNN idea but add gating mechanisms to better preserve long-range information and mitigate vanishing gradients.
    \item \textbf{Basic seq2seq has a severe bottleneck.} Expecting one final encoder vector to store the entire meaning of a sentence is ``wild,'' even if it works surprisingly well with enough data.
    \item \textbf{Attention is named well because it formalizes selective focus.} Instead of freezing all meaning into one vector, the decoder can revisit the source and decide what matters at each output step.
    \item \textbf{Attention is useful but not a perfect explanation of model behavior.} Attention heatmaps are informative, but they should not be mistaken for full causal interpretability.
\end{itemize}

\subsection{Examples}
\begin{itemize}
    \item \textbf{Long-range agreement example:} pronoun prediction in
    \[
    \text{``Emily earned the top grade \dots Everyone was proud of her''}
    \]
    versus
    \[
    \text{``Miguel earned the top grade \dots Everyone was proud of him''}
    \]
    requires preserving information from many tokens earlier.
    \item \textbf{Feed-forward limitation:} a fixed-window network predicting the next word in ``Anqi was late for \_\_\_'' can use nonlinearities but still lacks persistent memory.
    \item \textbf{RNN generation:} beginning with a start token, an RNN may generate text such as ``Sorry Harry shouted, panicking \dots'' one token at a time until an end token is produced.
    \item \textbf{Seq2Seq translation example:}
    \[
    \text{The brown dog ran} \rightarrow \text{Le chien brun a couru}
    \]
    where the output length differs from the input length.
    \item \textbf{Attention weighting example:} for the first decoder word, the model may assign weights such as \(0.51, 0.28, 0.14, 0.07\) over the encoder states, then form a weighted context vector.
    \item \textbf{Applications beyond translation:}
    \begin{itemize}
        \item many-to-one sentiment classification,
        \item one-to-one sequence tagging such as part-of-speech tagging,
        \item constituency parsing,
        \item image captioning with visual attention.
    \end{itemize}
\end{itemize}

\subsection{Common Pitfalls / Limitations}
\begin{itemize}
    \item \textbf{Feed-forward networks do not model sequential state.} Every prediction is independent except for the manually provided window.
    \item \textbf{RNNs are hard to parallelize.} Computation is sequential because each hidden state depends on the previous one.
    \item \textbf{BPTT is expensive.} Long sequences make training slow.
    \item \textbf{Vanishing gradients} make far-away context effectively disappear during learning.
    \item \textbf{Exploding gradients} destabilize optimization; gradient clipping helps numerical stability but is not a complete solution.
    \item \textbf{RNNs still struggle to use long-range context well in practice,} even though they are theoretically capable of representing it.
    \item \textbf{Vanilla seq2seq compresses the full source sequence into one vector,} which is a severe information bottleneck.
    \item \textbf{Attention context vectors can still ignore positional structure if built only as weighted sums.} This foreshadows the need for explicit positional mechanisms in later architectures.
    \item \textbf{Attention visualizations are not guaranteed explanations.} High attention to a word does not necessarily imply causal responsibility for a prediction.
\end{itemize}

\subsection{Connections to Other Lectures}
\begin{itemize}
    \item This lecture directly continues Lecture 2. Word2vec produced type-level embeddings with no contextualization; RNNs turn sequence modeling into a way to learn \textbf{contextualized} token representations.
    \item Lecture 2's log-linear and bag-of-words models ignored order or used only shallow local windows; RNNs add persistent hidden state and shared parameters across positions.
    \item Seq2seq extends the one-to-one and many-to-one sequence models from this lecture into \textbf{variable-length} input/output prediction, especially for translation.
    \item Attention is the bridge to the next lecture: once the decoder learns to compare states and reweight all source positions, the natural next step is to remove recurrence and build models around attention itself, i.e.\ Transformers.
    \item The historical machine-translation discussion also connects back to Lecture 2's Bayes-style factorization \(p(y \mid x) \propto p(x \mid y)p(y)\); seq2seq replaces the old modular pipeline with one learned end-to-end model.
\end{itemize}

\subsection{Exam-Relevant Summary}
\begin{itemize}
    \item \textbf{Word token vs.\ word type:} RNN hidden states are contextualized token representations; word2vec embeddings are type-based.
    \item \textbf{Feed-forward LM weakness:} fixed window, no memory, poor parameter efficiency as context grows.
    \item \textbf{Basic RNN equations:}
    \[
    h_t = f(Wx_t + Vh_{t-1} + b_h), \qquad
    \hat{y}_t = \operatorname{softmax}(Uh_t + b_y)
    \]
    \item \textbf{Initial state:}
    \[
    h_0 = \mathbf{0}
    \]
    \item \textbf{Training loss:}
    \[
    \mathcal{L} = \sum_{t=1}^{T} \operatorname{CE}(y_t,\hat{y}_t)
    \]
    \item \textbf{Teacher forcing:} use gold previous tokens during training.
    \item \textbf{BPTT:} gradients propagate through the unrolled sequence; this is expensive and causes vanishing/exploding gradient problems.
    \item \textbf{Truncated BPTT:} truncate the backward pass, not the recurrent forward memory.
    \item \textbf{Gradient clipping:} controls exploding gradients by limiting gradient norm.
    \item \textbf{RNN strengths:} shared weights across positions, stateful memory, variable-length sequential processing.
    \item \textbf{RNN weaknesses:} slow training, poor parallelism, vanishing/exploding gradients, limited practical use of very long context.
    \item \textbf{Seq2Seq:} encoder reads the source sequence, decoder generates the target sequence.
    \item \textbf{Basic seq2seq bottleneck:} all source meaning must pass through the final encoder hidden state.
    \item \textbf{Attention equations:}
    \[
    e_{t,i} = \operatorname{score}(h^{dec}_t, h^{enc}_i), \qquad
    a_{t,i} = \frac{\exp(e_{t,i})}{\sum_j \exp(e_{t,j})}, \qquad
    c_t = \sum_i a_{t,i} h^{enc}_i
    \]
    \item \textbf{Attention takeaway:} the decoder conditionally focuses on different encoder states at each output step, greatly improving seq2seq performance.
    \item \textbf{Historical significance:} RNN encoder-decoder models with attention were the major step before attention-only Transformer models.
\end{itemize}

\end{document}
